\documentclass[11pt,a4paper,twoside]{article}
\usepackage[utf8]{inputenc}
\usepackage[english]{babel}
\usepackage{actuarialangle, actuarialsymbol}
\usepackage{amsmath,amsfonts,amsthm,amssymb,mathtools}

\makeatletter
\newcommand{\ix}{\@actothersc{i}}
\makeatother

\begin{document}


Decompose the health insurance premium according to German regulations
and list its calculation bases and value categories.

+ Risk Premium
+ Savings Premium
= Savings Premium
+ Safety Margin
+ Surcharge Standard tariff (only optional tariffs)
+ Costs surcharge
= Gross premium
+ 10\% Statutory Gross Premium Surcharge Until Age 60

Calculation bases are
- Actuarial interest
- Elimination types
- Head claims
- Safety margin
- Surcharge for standard tariff in optional tariffs
- Costs surcharge

For the calculation bases, we distinguish
- true values: abstract mathematical values,
- effective values: current empirical values,
- calculatory values: values used in the calculation of tariffs, often including safety margins.



\iffalse
Let the total population at time \( t \) be  and \( t \in \text{KB} \), where  .
Define  as the set of ,
and 
Let the  stochastic time process be 
The true head claim for age x at time t.
\fi

Define head claims, profiles and effective calculatory head claims


Definitions
\( K \): head claim, expected value of insurance benefits in a certain time period for a single risk.
\( I(t) \): insurees at time \( t \), \( t \in \text{KB} \)
\( I = \cup_{t \in \text{KB}} I(t) \): insurees at time \( t \), \( t \in \text{KB} \)
\( \text{KB} \): Set of admissable dates (usually one calender year increments)
\( \text{AB} = \{x_{\text{min}}, \dots, x_{\text{max}}\} \): admissable ages of insurees
\( I_{x}(t) = \{ i \in I(t)| \text{age}(i) = x_i(t) = x \} \): individuals aged \( x \)
\( (s_i(t))_{t \in \text{KB}} \): stoch. process of medical expenses for individual \( i \in I \)
\( K_x(t) = \mathbb{E}[s_i(t)] \): true head claim for age \( x \) at time \( t \).

Profiles and their corresponding head claims are motivated by the:
Method of Rusam
The idea relies on a multiplicative splitting of tables of head claims in an age depending profile and a
normalization factor, the basic head claim.
\( G(t) = K_{x_0}(t) \) : true basic head claim at selected age \( x_0 \)
\( (k_x(t))_{x \in \text{AB}} = \left(\frac{K_x(t)}{G(t)} \right)_{x \in \text{AB}} \) : profile of \(  \)

\( K_x^{\text{eff}}(t_0) = \frac{S_x(t_0) - \text{RZ}_x(t_0) - \text{SO}_{x}(t_0)}{L_x(t_0)} \): effective head claims
\( S_x(t_0) \): effective claims expenditure for \( I_x(t_0) \)
\( \text{RZ}_x(t_0) \): sum of risk surcharges for \( I_x(t_0) \)
\( \text{SO}_{x}(t_0) \): expenditures due to special effects for \( I_x(t_0) \), e.g. pandemics, war
\( L_x(t_0) = |I_x(t_0)| \)

\( k_x^{\text{eff}}(t_0) = \frac{K_x^{\text{eff}}(t_0)}{K_{x_0}^{\text{eff}}(t_0)} \): effective profile at selected age \(x_0\)
\( G^{\text{eff}}(t_0) = \frac{S(t_0) - \text{RZ}(t_0) - \text{SO}(t_0)}{\sum_{x} L_x(t_0) \cdot k_x^{\text{eff}}(t_0)} \) : effective basic head claim
\( S(t_0) = \sum_{x} S_x(t_0) \) and analogously \( \text{RZ},\text{SO}  \).

Opposed to the effective head claims, the calculatory head claims are obtained via extrapolation:
1. Choose calculatory profile \( k_x^{\text{calc}} \) by smoothing of observed values or external data.
2. Estimation of basic head claims in the observation years
\[
 \hat{G}(t_j) = \frac{S(t_j) - \text{RZ}(t_j) - \text{SO}(t_j)}{\sum_{x} L_x(t_j) \cdot k_x^{\text{calc}}}
\]
for \( j = 0, -1, -2 \).
3. Choose calculatory head claims \( G^{\text{calc}}(t_1) \) by extrapolation of value from step 2.
4. Choose calculatory head claims via Rusam-ansatz by \( K_x^{\text{calc}}(t_1) = G^{\text{calc}}(t_1) k^{\text{calc}}_x \)

In case of a similar population structure, the profile is approximately independent of tariff and
company. Thus, profile data, published by BaFin or the PKV association, can be used to
implement new tariffs.
Over a period of few years, the profiles are approximately independent of time. Therefore, the
sequence of head claims is determined by the evolution of the basic head claim over time.
For empirical data gathering of head claims, the basic head claims and the profile are determined
separately.

Last the value of the net benefits is then given by
\[
  A_x(t) = \sum_{\nu = 0}^{x_{\text{max}} - x} v^\nu K_{x + \nu} {}_{\nu}p_x
\]
with \( _\nu p_x = \mathbb{P}[i \in I_{x + \nu}(t + \nu) | i \in I_x(t)]\).
and the corresponding net premium is \( P_x = \frac{A_x(t)}{\ddot{a}_x} \)



Define the aging reserve, premium splitting and the transfer value

The aging reserve is given by:
\[ 
  \Vx[m]{x} = \sum_{\nu = 0}^{x_{\text{max}} - (x + m)} v^\nu \cdot \px[\nu]{x + m} \cdot K_{x+m+\nu} - P_x \sum_{\nu = 0}^{x_{\text{max}} - (x + m)} v^\nu \cdot \px[\nu]{x+m} 
\]

The prospective aging Reserve on the basis of net annual premiums is computed as
\[ \Vx[m]{x} = (A_{x+m} - P_x)\ax**{x+m} \]

The contribution to the aging reserve can be decomposed into individual parts
\( \Vx[m+1]{x} - \Vx[m]{x} = (P_x - K_{x+m})(1 + i) + i_m V_x + (\qx{x+m} + w_{x+m}) \cdot \Vx[m+1]{x} \) \\
\( (P_x - K_{x+m})(1 + i) \): premium surplus with added accrued interest, \\
\( i_m V_x \): calculatory interest income on the aging reserve at the beginning of the year \\
\( (\qx{x+m} + w_{x+m}) \cdot \Vx[m+1]{x} \): calculatory inheritance of the aging reserve at the end of the year, due to exit from \( I_x \). \\

This also allows to split the premium

\[P_x = (v\cdot \Vx[m+1]{x} - \Vx[m]{x}) + K_{x+m} - v(\qx{x+m} + w_{x+m})\cdot \Vx[m+1]{x} \]
\( (v\cdot \Vx[m+1]{x} - \Vx[m]{x}) \): savings premium as the change in the aging reserve \\
\( K_{x+m} - v(\qx{x+m} + w_{x+m})\cdot \Vx[m+1]{x} \): risk premium, consisting of the natural premium \( K_{x+m} \)
and the premium reducing impact of the aging reserve \( v(\qx{x+m} + w_{x+m})\cdot \Vx[m+1]{x} \). \\


According to German law, the transfer value for an insurance contract must be calculated in the following terms: \\
\( \actsymb[][(\alpha)]{B}{x} \) : Annual gross premium at age \(x\) after \( m \) contract years  \\
\( \Vx[m]{x}[\text{GZ}] \): Aging reserve from the statutory surcharge \\
\( \Vx[m]{x}[g] \): Aging reserve from medical expenses tariff including ported benefits of the transfer value \\
\( \Vx[m]{x}[\text{BT}] \): (Fictitious) Aging reserve in the basic tariff if the person would have been insured in the 
basic tariff since contract inception. \\

Then the transfer value is given by
\[ \actsymb[m]{T}{x} = \Vx[m]{x}[\text{GZ}] + \max(0, \min(\Vx[m]{x}[g] + \max(0,  \actsymb[][(\alpha)]{Z}{x}), \Vx[m]{x}[\text{BT}])) \]
where \( \actsymb[][(\alpha)]{Z}{x} = \frac{5 - m}{5} \cdot \alpha \cdot \actsymb[][(\alpha)]{B}{x} \),
represents the evenly distributed zillmerization for the first five years.


Explain the old age premium problem

Causes for the premium problem of older insured persons
1. Increasing medical expenses over time
2. The mathematical problem of old age
3. Increasing profiles over time
4. Decreasing elimination rates over time
5. Lower income after start of retirement

Point 2 refers to the result that under the assumptions
- the elimination structure does not change at adjustment
- the aging reserve is strictly positive
a premium adjustment for insuree aged \( x \) after \( m \) contract years,
which is represented by an increase of \( 1 + j \) of the head claims, satisfies:

\[
  \frac{P_{x / x+m }}{P_x} = 1 + j + \frac{j}{P_x} \frac{{}_m V_x}{\ddot{a}_{x+m}} > 1 + j
\]
so the increase of the net premium is proportionally larger than the increase of the head claims.

The potential remedies of this so called old age premium problem are categorized
1. Positive influences on gross premiums in old age; however, benefits and premiums of
   equivalence equation for the net premium are unchanged.
2. Change the net premium and therefore the gross premium over the contract time
   such that old age premiums are reduced
3. Influences on the development of head claims or limits on insurance benefits.

Concrete Examples are
1. Using unit costs surcharges, excess interest, repurposing the reserve for premium refunds
2. dynamic tariffs, statutory 10\%-surcharge by law.
3. Lower insurance while qualitatively keep the most relevant insurance level (extended retention, lower daily allowances), 
   Change in standard / basic tariff



In health insurance, how is determined if the premia need adjustments?

According to §14 KalV the triggering value is the sole criterion for premia adjustments.
The triggering factor serves the purpose to compare required and calculatory benefits.
The factor is computed separately for different groups and tariffs.
The algorithm to compute it goes like this:

1. Compute with the last 3 periods \( t_0, t_{-1}, t_{-2} \) the calculatory basic head claims \( G^{\text{calc}}(t_{1}) \)
and associated factors.

2. Compute the last 3 effective basic head claims \( G^{\text{eff}}(t_j) \).

3. Extrapolate the effective basic head claims according to the formula (linear regression):
\[ G^{\text{req}}(t_{2}) = \frac{3}{2}(G^{\text{eff}}(t_0) - G^{\text{eff}}(t_{-2})) + \frac{1}{3} \sum_{i=-2}^0 G^{\text{eff}}(t_{i}) \]
4. The triggering factor is defined as \( \text{AF}(t_{1}) = \frac{G^{\text{req}}(t_{2})}{G^{\text{calc}}(t_{1})} \)
(actually its defined differently but one can derive this equality)

The criterion for a premium adjustments at time \( t_1 = t_0 + 1 \) is then 
\( |1 - \text{AF}(t_{1})| > 0,1 \)

Alternatively this might be approximated with cash values for required (empirical) and calculated mortality rates 
\( A^{\text{req}}_x, A^{\text{calc}}_x \):
Let \( \text{VQ}_x^{\text{mort}} = \frac{A^{\text{req}}_x}{A^{\text{calc}}_x} \) and take the averages 
over the age intervals \( [21, 45], [46, 70], [71, 95] \), denote them by 
\( \text{AVQ}_i^{\text{mort}} = \frac{1}{25} \sum_x \text{VQ}_x^{\text{mort}} \), \( i \in \{1, 2, 3\}\).

The triggering factor is then 
\( \text{AF}^{\text{mort}} = \max_i \text{AVQ}_i^{\text{mort}} \)
and the corresponding limit
\( |1 - \text{AF}(t_{1})| > 0,05 \)

Any other adjustments to premia can only be done if it is not to the disadvantage of the insured persons.
\iffalse
Compute \(S^{\text{req}}(t_{2}) = \sum_{x = x_{\text{min}}}^{x_{\text{max}}} K_x^{\text{req}}{x}(t_{2}) \) 
assuming that \( I(t_{2}) = I(t_{0}) \) and the same profile, thus \( K_x^{\text{req}}{x}(t_{2}) = G^{\text{req}}(t_{2}) k_x^\text{calc} \).
\fi


\end{document}
